% ============================================================
%  Informe Técnico — Práctica Final Integradora
%  Análisis y Diseño de Algoritmos · Universidad EAFIT · 2026-1
% ============================================================
\documentclass[12pt, a4paper]{article}

% ── Codificación e idioma ──────────────────────────────────
\usepackage[spanish, es-noshorthands]{babel}
\usepackage[utf8]{inputenc}
\usepackage[T1]{fontenc}
\usepackage{lmodern}
\usepackage{microtype}

% ── Geometría ─────────────────────────────────────────────
\usepackage[
  a4paper,
  top=2.5cm, bottom=2.5cm,
  left=3cm, right=2.5cm,
  headheight=14pt
]{geometry}

% ── Encabezado y pie ──────────────────────────────────────
\usepackage{fancyhdr}
\pagestyle{fancy}
\fancyhf{}
\fancyhead[L]{\small\textit{Práctica Final Integradora — ADA 2026-1}}
\fancyhead[R]{\small\textit{Universidad EAFIT}}
\fancyfoot[C]{\small\thepage}
\renewcommand{\headrulewidth}{0.4pt}

% ── Matemáticas ───────────────────────────────────────────
\usepackage{amsmath, amssymb}

% ── Tablas ────────────────────────────────────────────────
\usepackage{booktabs, tabularx, array}
\renewcommand{\arraystretch}{1.25}

% ── Listas ────────────────────────────────────────────────
\usepackage{enumitem}
\setlist{topsep=4pt, itemsep=2pt, parsep=0pt}

% ── Código ────────────────────────────────────────────────
\usepackage{listings}
\lstset{
  basicstyle=\small\ttfamily,
  frame=single,
  numbers=left,
  numberstyle=\tiny\color{gray},
  numbersep=6pt,
  breaklines=true,
  captionpos=b,
  xleftmargin=10pt,
  aboveskip=6pt,
  belowskip=6pt,
  showstringspaces=false,
}
\usepackage{xcolor}

% ── Captions e hipervínculos ──────────────────────────────
\usepackage{caption}
\captionsetup{font=small, labelfont=bf, skip=4pt}
\usepackage[colorlinks=true, linkcolor=black, urlcolor=black]{hyperref}

% ── Párrafos ──────────────────────────────────────────────
\setlength{\parindent}{0pt}
\setlength{\parskip}{6pt}

% ============================================================
%  PORTADA
% ============================================================
\begin{document}
\thispagestyle{empty}

\begin{center}
  \vspace*{1.5cm}
  {\large\bfseries Universidad EAFIT}\\[4pt]
  {\normalsize Departamento de Ciencias de la Computación}\\[2pt]
  {\normalsize Análisis y Diseño de Algoritmos · 2026-1}

  \vspace{2.5cm}
  \rule{\linewidth}{0.6pt}\\[10pt]
  {\LARGE\bfseries Práctica Final Integradora}\\[8pt]
  {\large Informe Técnico}\\[10pt]
  \rule{\linewidth}{0.6pt}

  \vspace{1cm}
  {\normalsize\itshape
    Optimización de rutas y planificación de recursos\\
    en redes de telecomunicaciones}

  \vspace{2cm}
  \begin{tabular}{rl}
    \textbf{Integrante 1:} & Nathalia Valentina Cardoza Azuaje \\[4pt]
    \textbf{Integrante 2:} & Andres Felipe Velez Alvarez \\[4pt]
    \textbf{Integrante 3:} & Sebastian Salazar Henao
  \end{tabular}

  \vspace{1.5cm}
  {\small Repositorio: \texttt{ADA\_PF\_Cardoza\_Salazar\_Velez}}\\[4pt]
  {\small 8 de mayo de 2026}
\end{center}

\newpage

% ── Tabla de contenidos ───────────────────────────────────
\tableofcontents
\newpage

% ============================================================
\section{Introducción}
% ============================================================
\textit{[Completar al final con el contexto general del problema y la justificación de por qué cada módulo requiere su paradigma específico.]}

% ============================================================
\section{Descripción del Dataset}
% ============================================================
\textit{[Completar con estadísticas del CSV cargado, manejo de valores nulos en \texttt{TotalCharges} y procedimiento de construcción del grafo.]}

% ============================================================
\section{Módulo A — Divide y Vencerás}
% ============================================================
\textit{[Pseudocódigo de MergeSort y búsqueda binaria; análisis de complejidad; tabla de tiempos empíricos; resultados de las 5 consultas.]}

% ============================================================
\section{Módulo B — Algoritmo Codicioso (Kruskal)}
% ============================================================
\textit{[Pseudocódigo de Kruskal con Union-Find; demostración del Lema del ciclo; lista de aristas del MST y peso total; verificación manual en subgrafo de 5 nodos.]}

% ============================================================
\section{Módulo C — Programación Dinámica (Mochila 0-1)}
% ============================================================

\subsection{Planteamiento}

Se seleccionan las \textbf{50 solicitudes activas} (\texttt{Churn = No}) de mayor prioridad (\texttt{tenure}), extraídas de la salida ordenada del Módulo A. Cada solicitud $i$ se modela como un ítem de la mochila:

\begin{itemize}
  \item \textbf{Peso} $w_i = \texttt{round(TotalCharges)}$: ancho de banda requerido (entero).
  \item \textbf{Valor} $v_i = \texttt{round(MonthlyCharges} \times 10)$: ingreso mensual en centavos.
\end{itemize}

El objetivo es encontrar el subconjunto $S$ que maximice $\sum_{i \in S} v_i$ sin superar la capacidad $W$, con variables de decisión binarias $x_i \in \{0, 1\}$. 

\item \textbf{Justificacion de escenarios.} En la entrega se especifica $W=500$; no obstante, los pesos encontrados en el top 50 de solicitudes activas por tenure son todos mayores a 500 (en unidades de TotalCharges redondeado). Al aplicar PD con $W=500$, la mochila resulta vacia porque ningun item cabe, y el valor optimo es trivialmente 0. Para observar el comportamiento real de ambos paradigmas (PD vs. codicioso), se decidio evaluar adicionalmente $W=5000$, una capacidad que permite seleccionar items significativos y revelar las diferencias en las estrategias de decision.

\subsection{Recurrencia y tabulación}

Se define $dp[i][w]$ como el valor máximo obtenible usando los primeros $i$ ítems con capacidad $w$:

\[
dp[i][w] =
\begin{cases}
  0, & \text{si } i = 0 \text{ o } w = 0 \\[3pt]
  dp[i-1][w], & \text{si } w_i > w \\[3pt]
  \max\!\bigl\{dp[i-1][w],\; dp[i-1][w - w_i] + v_i\bigr\}, & \text{si } w_i \leq w
\end{cases}
\]

La subestructura óptima garantiza que la solución global se construye correctamente a partir de subproblemas independientes.


\textbf{Pseudocódigo — tabulación.}
\begin{lstlisting}
for i = 1 to n:
    wi = items[i-1].peso;  vi = items[i-1].valor
    for w = 0 to W:
        if wi > w:
            dp[i][w] = dp[i-1][w]
        else:
            dp[i][w] = max(dp[i-1][w], dp[i-1][w - wi] + vi)
\end{lstlisting}

\subsection{Reconstrucción (backtracking)}

Para recuperar el conjunto exacto de ítems seleccionados se recorre la tabla desde $(n, W)$ hacia $(1, 0)$:

\begin{lstlisting}
w = W
for i = n downto 1:
    if items[i-1].peso <= w  and  dp[i][w] != dp[i-1][w]:
        seleccionar items[i-1]
        w = w - items[i-1].peso
\end{lstlisting}

Si $dp[i][w] \neq dp[i-1][w]$, el ítem $i$ fue incluido en la solución óptima; de lo contrario, fue excluido.

\subsection{Resultados obtenidos}

\subsubsection*{Escenario $W = 500$}

\begin{itemize}
  \item Valor óptimo: \textbf{0} centavos. Solicitudes seleccionadas: \textbf{ninguna}.
\end{itemize}

Este resultado es correcto: todos los pesos $w_i = \texttt{round(TotalCharges)}$ de las 50 solicitudes activas de mayor prioridad son estrictamente superiores a 500, por lo que ningún ítem cabe en la mochila. La escala real de \texttt{TotalCharges} para clientes con alta antigüedad supera ampliamente las 500 unidades.

\subsubsection*{Escenario $W = 5\,000$}

\begin{itemize}
  \item Valor óptimo total: \textbf{707} centavos.
  \item Solicitudes seleccionadas: \texttt{9286-BHDQG} y \texttt{4312-KFRXN}.
  \item Capacidad utilizada: \textbf{4\,851} / 5\,000 unidades.
\end{itemize}

\subsection{Fallo del enfoque codicioso}

La heurística codiciosa por ratio $v_i/w_i$ selecciona ítems en orden decreciente de rentabilidad hasta agotar la capacidad. Para demostrar que esta estrategia no garantiza optimalidad, se construyó un contraejemplo explícito con exactamente 3 solicitudes del conjunto anterior:

\begin{table}[h!]
\centering
\small
\begin{tabularx}{\textwidth}{
  >{\raggedright\arraybackslash}p{4cm}
  >{\raggedright\arraybackslash}X
  >{\centering\arraybackslash}p{2cm}
  >{\centering\arraybackslash}p{1.8cm}
}
\toprule
\textbf{Enfoque} & \textbf{Solicitudes seleccionadas} & \textbf{Valor total} & \textbf{¿Óptimo?} \\
\midrule
Codicioso (ratio $v/w$) & \texttt{4312-KFRXN}, \texttt{6734-PSBAW}, \texttt{7606-BPHHN} & 688 & No \\
PD (Mochila 0-1) & \texttt{9286-BHDQG}, \texttt{4312-KFRXN} & \textbf{707} & Sí \\
\bottomrule
\end{tabularx}
\caption{Contraejemplo: heurística codiciosa \textit{vs.} solución óptima por PD ($W = 5\,000$).}
\end{table}

El codicioso selecciona 3 ítems con peso total 4\,904 y valor 688 siguiendo el criterio local. La PD descarta uno de ellos e incluye \texttt{9286-BHDQG}, logrando un valor de 707 con solo 2 ítems. Esto evidencia que el criterio local óptimo (mayor ratio) puede bloquear combinaciones globalmente superiores cuando los pesos son heterogéneos.

\subsection{Complejidad y pseudopolinomialidad}

\paragraph{Tiempo.} La tabulación recorre $(n+1) \times (W+1)$ celdas, cada una en $O(1)$:
\[
  T(n, W) = \Theta(n \cdot W).
\]
El backtracking aporta $O(n)$ adicional, dominado por el término anterior.

\paragraph{Espacio.} La tabla completa requiere $\Theta(n \cdot W)$ de memoria. Si solo se necesita el valor óptimo sin reconstrucción, se puede reducir a $\Theta(W)$ conservando únicamente dos filas consecutivas.

\paragraph{Pseudopolinomialidad.} El costo $\Theta(n \cdot W)$ es polinomial en el valor numérico de $W$, pero $W$ se codifica en $\lceil \log_2 W \rceil$ bits; por tanto, el algoritmo es exponencial en el tamaño de la entrada en bits. Esto al algoritmo como \textbf{pseudopolinomial}. En este proyecto, el redondeo de pesos y el factor $\times 10$ en valores son necesarios para trabajar con índices enteros, pero incrementan $W$ y, con ello, el costo computacional.

\subsection{Discusión}

\begin{itemize}
  \item \textbf{Escenario $W = 500$.} El valor óptimo $0$ no es un error de implementación, sino consecuencia directa de la escala del dataset: los clientes de mayor antigüedad acumulan cargos totales muy superiores a 500.

  \item \textbf{Escenario $W = 5\,000$.} La PD supera al codicioso en 19 centavos, confirmando que la exploración global del espacio de combinaciones es necesaria cuando los pesos no son proporcionales a los valores.

  \item \textbf{Garantía de optimalidad.} A diferencia del enfoque codicioso, la PD garantiza encontrar la solución exacta, a costa de un costo que crece linealmente con $W$.
\end{itemize}

% ============================================================
\section{Integración del Pipeline}
% ============================================================
\textit{[Diagrama de flujo del programa completo; descripción de cómo la salida del Módulo A alimenta al Módulo C.]}

% ============================================================
\section{Conclusiones}
% ============================================================
\textit{[Comparación de los tres paradigmas en términos de garantías de optimalidad y complejidad; reflexión sobre el comportamiento en datos reales.]}

% ============================================================
\section{Herramientas Utilizadas}
% ============================================================
\textit{[Citación explícita de cualquier herramienta de IA generativa usada, conforme a las políticas del curso.]}

% ============================================================
\section{Referencias}
% ============================================================

\begin{enumerate}[label={[\arabic*]}]
  \item Blastchar. \textit{Telco Customer Churn}. Kaggle, 2018.
        \url{https://www.kaggle.com/datasets/blastchar/telco-customer-churn}
  \item Cormen, T.\,H., Leiserson, C.\,E., Rivest, R.\,L., \& Stein, C. (2022).
        \textit{Introduction to Algorithms} (4.ª ed.). MIT Press.
  \item Material del curso de Análisis y Diseño de Algoritmos — Universidad EAFIT, 2026-1.
\end{enumerate}

% ============================================================
\section{Roles del Equipo}
% ============================================================
\textit{[Contribución específica de cada integrante por módulo.]}

\end{document}
